\section{Introduction}
\label{sec:introduction}

As inverter-based resources (IBRs) displace synchronous generation, the strength of the grid at
their connection points becomes a planning concern \cite{SvK2025, fingrid2023, Denholm2022, NERC2018}. Fault-current-based metrics are the
usual first screen: they are quick to compute from a network model, and their values decide
which connections and operating conditions receive detailed electromagnetic-transient study
\cite{ieee2800,nerc2017,cigre671}. IEEE Std 2800-2022 adds that a value used across a system
should cover a variety of operating conditions and topologies \cite{ieee2800}. Screening is
therefore only as reliable as the metric's ability to identify when, and under which outages, a
site is weakest.

Several such metrics are in use, and each comes with its own conventions. The short-circuit ratio
(SCR) and the weighted SCR (WSCR) are evaluated at the point of interconnection on the network
without the IBR plants \cite{ieee2800,zhang2014}. The equivalent SCR (ESCR), titled SCRIF in IEEE
Std 2800-2022 \cite{ieee2800,cigre671}, adds the plants' interaction through multi-infeed interaction
factors \cite{cigre364}, and is applied on the medium-voltage side \cite{fingrid2023}. The
site-dependent SCR (SDSCR) weights the other plants through the impedance matrix at the converter
terminals \cite{wu2018}. Names and the construction of the interaction factor differ between
sources \cite{joswigjones2026,goonasekara2023,darii2026}. Comparisons of the metrics so far rest
on discrete cases: Wu et al. \cite{wu2018} compare them over twelve operating and
contingency conditions, others use a handful of loading cases or single snapshots
\cite{sanni2022,dhungana2026,khamees2026,rosengren2025}, and commitment has been shown to drive
WSCR within a single period \cite{kim2025}. WSCR, however, is a group metric and does not give the strength of individual sites.
To the authors' knowledge, no study has followed the
metrics through a time-ordered sequence of operating states, in which commitment changes with
load, together with a full set of line and transformer outages.

This paper asks whether SCR, WSCR, ESCR and SDSCR, each applied as specified, lead to consistent
screening conclusions along such a sequence, and what explains any difference. Over time, it asks
whether the metrics agree on when each site is weakest, whether any metric tells more than the
committed set of synchronous generators, and whether the usual planning snapshots find the weekly
minimum. Under outages, it asks what a single branch outage does to each site. Finally, it asks whether
differences come from the metrics themselves or from their conventions (network model, evaluation
point and interaction-factor construction).

The study adds three IBR plants to the 37-bus network of \cite{gloverPSAD} and evaluates the four metrics over
a 168-hour week, with load scaled by a public demand profile, under a largest-first commitment rule,
intact and for each single line and transformer outage, and with every unit online on the intact
network. The contributions are:
\begin{itemize}
  \item a time-ordered, per-site comparison of the four metrics over a week and the full N-1 branch
        set, each under its practitioner-specified network, evaluation point and denominator;
  \item an attribution of every agreement and disagreement to a term of the metric, confirmed by
        removing that term or holding it constant;
  \item a reconciliation of the metrics' conventions, including how the choice of interaction
        factor changes ESCR's response to outages;
  \item a reproducible implementation on public data and open code.
\end{itemize}

Section~\ref{sec:metrics} defines the metrics and their conventions, and Section~\ref{sec:study}
the benchmark, operating states and experiments. Section~\ref{sec:results} presents the results,
Section~\ref{sec:discussion} discusses their implications and limitations, and
Section~\ref{sec:conclusion} concludes.
